%!TEX root = ../main.tex
% APÉNDICE C — Detalles técnicos e implementación

\chapter{Detalles técnicos e implementación}\label{apx:tecnicos}

\section{Configuración del motor de recomendación}

El motor de recomendación es de reglas ponderadas: no utiliza un modelo
entrenado, por lo que no tiene hiperparámetros ni conjuntos de datos de
entrenamiento. Para cada salón o proveedor calcula componentes entre 0 y
1, los combina con los pesos de la tabla~\ref{tab:pesos-recomendacion}
y expresa el resultado como un puntaje de 0 a 100. Los pesos se definen
en un archivo de configuración y pueden ajustarse sin modificar el
código.

\begin{table}[H]
    \centering
    \caption{Pesos del puntaje de recomendación}
    \label{tab:pesos-recomendacion}
    \small
    \begin{tabular}{l l r}
        \toprule
        \tb{Se recomienda} & \tb{Componente} & \tb{Peso} \\
        \midrule
        Salones & Presupuesto & 0,40 \\
                & Capacidad   & 0,30 \\
                & Reseñas     & 0,30 \\
        \midrule
        Proveedores de servicios & Presupuesto & 0,35 \\
                                 & Ubicación   & 0,30 \\
                                 & Reseñas     & 0,35 \\
        \bottomrule
    \end{tabular}
\end{table}

\begin{table}[H]
    \centering
    \caption{Cálculo de cada componente del puntaje}
    \label{tab:componentes-recomendacion}
    \small
    \begin{tabularx}{\textwidth}{l X}
        \toprule
        \tb{Componente} & \tb{Regla} \\
        \midrule
        Presupuesto & 1 si el presupuesto del evento cubre el mínimo del
                      rango de precios; si no, disminuye en proporción a
                      cuánto excede ese mínimo al presupuesto. \\
        Capacidad   & 1 si la cantidad de invitados está dentro de la
                      capacidad del salón; 0 si supera el máximo; si es
                      menor que el mínimo, es la proporción entre ambos. \\
        Ubicación   & 1 si la zona de cobertura del proveedor coincide
                      con la ubicación del evento; 0,4 en caso contrario. \\
        Reseñas     & Promedio de puntuaciones de 1 a 5 llevado a la
                      escala de 0 a 1. \\
        \bottomrule
    \end{tabularx}
\end{table}

Cuando falta información de alguna de las dos partes (por ejemplo, un
salón sin capacidad cargada o un evento sin presupuesto), el componente
toma un valor neutro de 0,5, de modo que la falta de datos no penaliza
la sugerencia. Cada recomendación incluye además un motivo en lenguaje
natural construido a partir de los componentes que más influyeron.

\section{Datos utilizados}

El sistema no utiliza conjuntos de datos externos. El recomendador y el
simulador operan sobre los datos cargados en la propia plataforma:
salones, proveedores, reseñas, reservas y eventos.

\section{Arquitectura de la aplicación}

\subsection{Stack tecnológico}

\begin{lstlisting}[caption={Dependencias principales}, label={lst:dependencias}]
Servidor
  PHP                       ^8.2
  Laravel                   ^12.0
  Laravel Sanctum           ^4.3    (tokens y roles)
  PHPUnit                   ^11.5   (pruebas automatizadas)
  Laravel Pint, Mockery             (formato de codigo y mocks)
  Base de datos             SQLite (desarrollo) / MySQL (produccion)
                            El mismo codigo corre en ambas.
  Correo y cliente HTTP     componentes propios de Laravel
  Mercado Pago              Checkout Pro, solo modo de prueba (sandbox),
                            detras de la interfaz PasarelaDePago

Cliente
  React                     ^19.1
  TypeScript                ^4.9
  Material UI (MUI)         ^7.3 con Emotion
  React Router              ^7.8
  Axios                     ^1.11   (cliente HTTP)
  Leaflet / React-Leaflet   ^1.9 / ^5.0  (mapas)
  dnd-kit                           (arrastrar y soltar)
  jsPDF / jspdf-autotable           (generacion de PDF)
  Create React App          react-scripts 5
  Jest y Testing Library            (pruebas del cliente)

Herramientas
  Git y GitHub
  GitHub Actions            CI con tres tareas: servidor con SQLite,
                            servidor con MySQL y cliente (tipos y build)
  Node                      22
\end{lstlisting}

\subsection{Esquema de base de datos}

Se presentan las tablas centrales del marketplace, simplificadas. Cada
salón y cada proveedor se asocian a un usuario con rol propio, lo que
permite aislar los datos de cada negocio.

\begin{lstlisting}[language=SQL, caption={Tablas centrales del marketplace (extracto simplificado)}, label={lst:esquema}]
CREATE TABLE salones (
    id               BIGINT PRIMARY KEY AUTO_INCREMENT,
    user_id          BIGINT NULL REFERENCES users(id),
    nombre_comercial VARCHAR(255) NOT NULL,
    descripcion      TEXT NULL,
    ubicacion        VARCHAR(255) NULL,
    capacidad_min    INT NULL,
    capacidad_max    INT NULL,
    rango_precio_min DECIMAL(12,2) NULL,
    rango_precio_max DECIMAL(12,2) NULL,
    activo           BOOLEAN NOT NULL DEFAULT TRUE,
    verificado       BOOLEAN NOT NULL DEFAULT FALSE
);

CREATE TABLE proveedores_servicios (
    id               BIGINT PRIMARY KEY AUTO_INCREMENT,
    user_id          BIGINT NULL REFERENCES users(id),
    nombre_comercial VARCHAR(255) NOT NULL,
    categoria        ENUM('catering','fotografia','decoracion',
                          'sonido_iluminacion','entretenimiento') NOT NULL,
    zona_cobertura   VARCHAR(255) NULL,
    rango_precio_min DECIMAL(12,2) NULL,
    rango_precio_max DECIMAL(12,2) NULL,
    activo           BOOLEAN NOT NULL DEFAULT TRUE,
    verificado       BOOLEAN NOT NULL DEFAULT FALSE
);

CREATE TABLE reservas (
    id                    BIGINT PRIMARY KEY AUTO_INCREMENT,
    evento_id             BIGINT NOT NULL REFERENCES eventos(id),
    salon_id              BIGINT NULL REFERENCES salones(id),
    proveedor_servicio_id BIGINT NULL REFERENCES proveedores_servicios(id),
    estado                ENUM('solicitado','confirmado','cancelado',
                               'completado') NOT NULL DEFAULT 'solicitado',
    monto                 DECIMAL(12,2) NULL,
    fecha_solicitud       DATE NULL
);

CREATE TABLE resenas (
    id          BIGINT PRIMARY KEY AUTO_INCREMENT,
    reserva_id  BIGINT NOT NULL UNIQUE REFERENCES reservas(id),
    puntuacion  TINYINT NOT NULL,
    comentario  TEXT NULL
);

CREATE TABLE simulaciones_evento (
    id                  BIGINT PRIMARY KEY AUTO_INCREMENT,
    evento_id           BIGINT NOT NULL REFERENCES eventos(id),
    variables_entrada   JSON NOT NULL,
    variables_resultado JSON NOT NULL,
    aplicado            BOOLEAN NOT NULL DEFAULT FALSE
);
\end{lstlisting}

\subsection{Variables de entorno}

La configuración del sistema se realiza mediante variables de entorno.
Se listan los nombres y su propósito; los valores secretos nunca se
incluyen en el repositorio ni en este documento.

\begin{table}[H]
    \centering
    \caption{Variables de entorno principales}
    \label{tab:variables-entorno}
    \small
    \begin{tabularx}{\textwidth}{p{4.1cm} X p{4.1cm}}
        \toprule
        \tb{Variable} & \tb{Propósito} & \tb{Ejemplo de valor} \\
        \midrule
        \texttt{APP\_ENV}   & Entorno de ejecución del servidor & \texttt{production} \\
        \texttt{APP\_KEY}   & Clave de cifrado de la aplicación (secreta) & (no se incluye) \\
        \texttt{APP\_DEBUG} & Muestra detalles de errores; debe estar desactivada en producción & \texttt{false} \\
        \texttt{APP\_URL}   & Dirección pública del servidor & \texttt{https://api.ejemplo.com} \\
        \texttt{DB\_CONNECTION} & Tipo de base de datos & \texttt{mysql} \\
        \texttt{DB\_HOST}, \texttt{DB\_PORT}, \texttt{DB\_DATABASE}, \texttt{DB\_USERNAME} & Conexión a la base de datos & (según el entorno) \\
        \texttt{DB\_PASSWORD} & Contraseña de la base de datos (secreta) & (no se incluye) \\
        \texttt{FRONTEND\_URL} & Dirección del cliente web, usada en enlaces enviados por el servidor & \texttt{https://app.ejemplo.com} \\
        \texttt{CORS\_ALLOWED\_ORIGINS} & Orígenes autorizados a consumir la \gls{api} & \texttt{https://app.ejemplo.com} \\
        \texttt{MERCADOPAGO\_ACCESS\_TOKEN} & Credencial de la cuenta vendedora de prueba de Mercado Pago (secreta) & (no se incluye) \\
        \texttt{MERCADOPAGO\_WEBHOOK\_SECRET} & Clave para validar los avisos de pago de Mercado Pago (secreta) & (no se incluye) \\
        \texttt{PAGOS\_PORCENTAJE\_SENA} & Porcentaje del monto de la reserva que se cobra como seña & \texttt{30} \\
        \texttt{REACT\_APP\_API\_URL} & Dirección de la \gls{api} que usa el cliente web & \texttt{https://api.ejemplo.com/api} \\
        \bottomrule
    \end{tabularx}
\end{table}
